\chapter*{Abstract}
\addcontentsline{toc}{chapter}{Abstract}

This monograph presents a formal mathematical framework for analyzing the Indus Valley (Harappan) script through the paradigm of \textbf{operational mathematics} and \textbf{state-based computation}. Rather than approaching ancient symbols as fixed phonetic or semantic units, we model them as \textbf{state operators} that transform computational contexts through formal compositional rules.

\section*{Core Contribution}

We develop an axiomatic system consisting of five axioms (A1--A5) that govern symbol composition and state transformation:

\begin{itemize}
    \item \textbf{A1 (Class Primacy)}: Symbols belong to disjoint semantic classes defined by operational behavior
    \item \textbf{A2 (No New Classes)}: Composition preserves class membership
    \item \textbf{A3 (Restricted Composition)}: Only compatible classes can compose
    \item \textbf{A4 (Terminal Closure)}: Terminal symbols enforce sequence boundaries
    \item \textbf{A5 (Hierarchical Superposition)}: Modifiers apply in well-defined layers
\end{itemize}

From these axioms, we prove the \textbf{\TCOP{} theorem} (Total Closure of Operations Principle): Every valid symbol sequence represents a well-defined state transformation, and compositions are closed under the system's rules.

\section*{Comparative Analysis}

We establish deep structural parallels with ancient Chinese character transformations, particularly the radical \ren{} (\textit{rén}, ``person'') and its systematic modifications. This comparison provides empirical validation from a documented operational system.

\section*{Formal Verification}

The framework has been formally verified in four independent proof assistants:
\begin{itemize}
    \item Lean 4
    \item Coq
    \item Isabelle/HOL
    \item Agda
\end{itemize}

All axioms and the \TCOP{} theorem have been machine-checked for logical consistency.

\section*{Implementation}

We provide:
\begin{itemize}
    \item Complete Python implementation for computational analysis
    \item C++ implementation for performance-critical applications
    \item Reproducible data processing pipeline
    \item Statistical analysis of the first 44 Harappan symbols
\end{itemize}

\section*{Explicit Non-Claims}

This work \textbf{does NOT}:
\begin{itemize}
    \item Claim to decipher the Indus script phonetically
    \item Provide word-by-word translations of inscriptions
    \item Establish historical meanings of specific symbols
    \item Propose phonetic values for any signs
    \item Claim to ``solve'' or ``crack'' the Indus script
\end{itemize}

\section*{Epistemological Stance}

We study the \textbf{mathematics of symbol composition}, not historical linguistics. Our claims concern \textbf{operational structure}, not \textbf{semantic content}. The framework is agnostic about phonetic values, linguistic affiliations, and historical functions of the script.

\section*{Significance}

This work demonstrates that:
\begin{enumerate}
    \item Ancient symbol systems can be rigorously analyzed using formal methods
    \item State-based operational semantics provides a general framework beyond specific scripts
    \item Structural validation is possible without phonetic decipherment
    \item Cross-cultural comparison reveals deep operational patterns
\end{enumerate}

\section*{Keywords}

Indus Valley Script, Harappan Script, Operational Semantics, State-Based Computation, Formal Methods, Ancient Chinese Characters, Comparative Paleography, Proof Assistants, Symbol Systems, Non-Phonetic Analysis

\vfill

\noindent\textbf{Citation:} Draconis (2026). \textit{Indus Operational Mathematics: A Formal Framework for State-Based Symbol Systems}. \url{https://github.com/ShantiDraconis/indus-operational-mathematics}
